%
% This file is part of the project of
% National Cheng Kung University (NCKU) Thesis/Dissertation Template in LaTex.
% This project is hold at
%     <https://github.com/wengan-li/ncku-thesis-template-latex>
% by Wen-Gan Li.
%
% This project is distributed in the hope of usefuling to someone,
% you can redistribute it and/or modify it under the terms of the
% Attribution-NonCommercial-ShareAlike 4.0 International.
%
% You should have received a copy of the
% Attribution-NonCommercial-ShareAlike 4.0 International
% along with this project.
% If not, see <http://creativecommons.org/licenses/by-nc-sa/4.0/legalcode.txt>.
%
% Please feel free to fork it, modify it, and try it.
% Have fun !!!
%

% ----------------------------------------------------------------------------

% Reference from:
%   <https://www.sharelatex.com/learn/Theorems_and_proofs>

%    Definition       (定義)
%    Condition        (條件)
%    Theorem          (定理)
%    Lemma            (引理)
%    Example          (例子)
%    Corollary        (推論)
%    Proposition      (主張)
%    Remark           (備註)
%    Proof            (證明)
%    Conjecture       (猜想)
%    Note             (附註)
%    Annotation       (註解)
%    Claim            (主張)
%    Case             (情況)
%    Acknowledgment   (確認)
%    Conclusion       (結論)
%    Criterion        (標準)
%    Assertion        (斷言)
%    Problem          (問題)
%    Question         (問題)
%    Hypothesis       (假設)
%    Summary          (總結)

% ----------------------------------------------------------------------------

% Shared functions

\pgfkeys
{
  /InsertTheoremOptions/.is family, /InsertTheoremOptions,
  default/.style =
  {
    title = \empty,
    label = \empty,
  },
  title/.estore in = \TmpValueTitle,
  label/.estore in = \TmpValueLabel,
} % End of \pgfkeys{}

\newcommand{\SetTheoremContentLabel}[1][\empty]
{%
  \ifthenelse{\equal{#1}{\empty}}{}{\IfNoValueF{#1}{\label{#1}}}%
} % End of \DeclareDocumentCommand{}

\newcommand{\InsertTheoremContent}[3][\empty]
{%
  % Parse the input
  \pgfkeys{/InsertTheoremOptions, default, #1}%
  %
  \ifthenelse{\equal{\TmpValueTitle}{\empty}}%
  {%
    \begin{#2}%
  }%
  {%
    \begin{#2}[\TmpValueTitle]%
  }%
  % Freeze the optional title before \label writes nameref metadata. Without this
  % copy the auxiliary file stores the mutable \TmpValueTitle control sequence,
  % which is reset by the next insertion.
  \ifthenelse{\equal{\TmpValueTitle}{\empty}}{}%
    {\expandafter\let\csname @currentlabelname\endcsname\TmpValueTitle}%
  %
  \SetTheoremContentLabel[\TmpValueLabel]%
  %
  #3%
  %
  \ifthenelse{\equal{#2}{\GetTheoremProofFormatEnvironmentName}}%
    {$\quad\blacksquare$}{}%
  %
  \end{#2}%
} % End of \newcommand{}

\newcommand{\MappingTheoremCounterResetCounter}[2][\empty]
{%
  \ifthenelse{\equal{#1}{Definition}}%
  {%
    \renewcommand{\GetTheoremDefinitionFormatFollowCounter}{#2}%
  }{}%
  %
  \ifthenelse{\equal{#1}{Condition}}%
  {%
    \renewcommand{\GetTheoremConditionFormatFollowCounter}{#2}%
  }{}%
  %
  \ifthenelse{\equal{#1}{Theorem}}%
  {%
    \renewcommand{\GetTheoremTheoremFormatFollowCounter}{#2}%
  }{}%
  %
  \ifthenelse{\equal{#1}{Lemma}}%
  {%
    \renewcommand{\GetTheoremLemmaFormatFollowCounter}{#2}%
  }{}%
  %
  \ifthenelse{\equal{#1}{Example}}%
  {%
    \renewcommand{\GetTheoremExampleFormatFollowCounter}{#2}%
  }{}%
  %
  \ifthenelse{\equal{#1}{Problem}}%
  {%
    \renewcommand{\GetTheoremProblemFormatFollowCounter}{#2}%
  }{}%
  %
  \ifthenelse{\equal{#1}{Corollary}}%
  {%
    \renewcommand{\GetTheoremCorollaryFormatFollowCounter}{#2}%
  }{}%
  %
  \ifthenelse{\equal{#1}{Proposition}}%
  {%
    \renewcommand{\GetTheoremPropositionFormatFollowCounter}{#2}%
  }{}%
  %
  \ifthenelse{\equal{#1}{Conjecture}}%
  {%
    \renewcommand{\GetTheoremConjectureFormatFollowCounter}{#2}%
  }{}%
  %
  \ifthenelse{\equal{#1}{Criterion}}%
  {%
    \renewcommand{\GetTheoremCriterionFormatFollowCounter}{#2}%
  }{}%
  %
  \ifthenelse{\equal{#1}{Assertion}}%
  {%
    \renewcommand{\GetTheoremAssertionFormatFollowCounter}{#2}%
  }{}%
  %
  \ifthenelse{\equal{#1}{Question}}%
  {%
    \renewcommand{\GetTheoremQuestionFormatFollowCounter}{#2}%
  }{}%
  %
  \ifthenelse{\equal{#1}{Hypothesis}}%
  {%
    \renewcommand{\GetTheoremHypothesisFormatFollowCounter}{#2}%
  }{}%
  %
} % End of \newcommand{}

\newcommand{\MappingTheoremCounter}[2][\empty]
{%
  \ifthenelse{\equal{#2}{Section}}%
  {%
    \MappingTheoremCounterResetCounter[#1]{section}%
  }{}%
  %
  \ifthenelse{\equal{#2}{Definition}}%
  {%
    \MappingTheoremCounterResetCounter[#1]{%
      \GetTheoremDefinitionFormatFollowCounter}%
  }{}%
  %
  \ifthenelse{\equal{#2}{Condition}}%
  {%
    \MappingTheoremCounterResetCounter[#1]{%
      \GetTheoremConditionFormatFollowCounter}%
  }{}%
  %
  \ifthenelse{\equal{#2}{Theorem}}%
  {%
    \MappingTheoremCounterResetCounter[#1]{%
      \GetTheoremTheoremFormatFollowCounter}%
  }{}%
  %
  \ifthenelse{\equal{#2}{Lemma}}%
  {%
    \MappingTheoremCounterResetCounter[#1]{%
      \GetTheoremLemmaFormatFollowCounter}%
  }{}%
  %
  \ifthenelse{\equal{#2}{Example}}%
  {%
    \MappingTheoremCounterResetCounter[#1]{%
      \GetTheoremExampleFormatFollowCounter}%
  }{}%
  %
  \ifthenelse{\equal{#2}{Proposition}}%
  {%
    \MappingTheoremCounterResetCounter[#1]{%
      \GetTheoremPropositionFormatFollowCounter}%
  }{}%
  %
  \ifthenelse{\equal{#2}{Conjecture}}%
  {%
    \MappingTheoremCounterResetCounter[#1]{%
      \GetTheoremConjectureFormatFollowCounter}%
  }{}%
  %
  \ifthenelse{\equal{#2}{Criterion}}%
  {%
    \MappingTheoremCounterResetCounter[#1]{%
      \GetTheoremCriterionFormatFollowCounter}%
  }{}%
  %
  \ifthenelse{\equal{#2}{Assertion}}%
  {%
    \MappingTheoremCounterResetCounter[#1]{%
      \GetTheoremAssertionFormatFollowCounter}%
  }{}%
  %
  \ifthenelse{\equal{#2}{Question}}%
  {%
    \MappingTheoremCounterResetCounter[#1]{%
      \GetTheoremQuestionFormatFollowCounter}%
  }{}%
  %
  \ifthenelse{\equal{#2}{Hypothesis}}%
  {%
    \MappingTheoremCounterResetCounter[#1]{%
      \GetTheoremHypothesisFormatFollowCounter}%
  }{}%
  %
  \ifthenelse{\equal{#2}{Problem}}%
  {%
    \MappingTheoremCounterResetCounter[#1]{%
      \GetTheoremProblemFormatFollowCounter}%
  }{}%
  %
  \ifthenelse{\equal{#2}{Corollary}}%
  {%
    \MappingTheoremCounterResetCounter[#1]{%
      \GetTheoremCorollaryFormatFollowCounter}%
  }{}%
  %
} % End of \newcommand{}

% ---------------------------------------------------------

% Resolve the configured parent counter before choosing the numbered or
% unscoped \newtheorem form. A followed type may itself be unscoped.
\newcommand{\NCKUPrivateInitNumberedTheoremFormat}[4]
{%
  \ifthenelse{\equal{#4}{\empty}}%
  {%
    \newtheorem{#2}{#3}%
  }%
  {%
    \MappingTheoremCounter[#1]{#4}%
    \ifthenelse{\equal{#4}{\empty}}%
    {%
      \newtheorem{#2}{#3}%
    }%
    {%
      \newtheorem{#2}{#3}[#4]%
    }%
  }%
} % End of \newcommand{}

% ---------------------------------------------------------

\pgfkeys
{
  /TheoremTheoremFormat/.is family, /TheoremTheoremFormat,
  default/.style =
  {
    EnvironmentName = {EnvTheorem},
    ShowText = {Theorem},
    FollowCounter = Section,
  },
  EnvironmentName/.estore in = \GetTheoremTheoremFormatEnvironmentName,
  ShowText/.estore in = \GetTheoremTheoremFormatShowText,
  FollowCounter/.estore in = \GetTheoremTheoremFormatFollowCounter,
} % End of \pgfkeys{}

\newcommand{\InsertTheorem}[2][\empty]
{%
  \InsertTheoremContent[#1]{\GetTheoremTheoremFormatEnvironmentName}{#2}%
} % End of \newcommand{}

\newcommand{\InitTheoremTheoremFormat}
{%
  \theoremstyle{plain}%
  \NCKUPrivateInitNumberedTheoremFormat{Theorem}{%
    \GetTheoremTheoremFormatEnvironmentName}{%
    \GetTheoremTheoremFormatShowText}{%
    \GetTheoremTheoremFormatFollowCounter}%
} % End of \newcommand{}

% ---------------------------------------------------------

\pgfkeys
{
  /TheoremDefinitionFormat/.is family, /TheoremDefinitionFormat,
  default/.style =
  {
    EnvironmentName = {EnvDefinition},
    ShowText = {Definition},
    FollowCounter = Section,
  },
  EnvironmentName/.estore in = \GetTheoremDefinitionFormatEnvironmentName,
  ShowText/.estore in = \GetTheoremDefinitionFormatShowText,
  FollowCounter/.estore in = \GetTheoremDefinitionFormatFollowCounter,
} % End of \pgfkeys{}

\newcommand{\InsertDefinition}[2][\empty]
{%
  \InsertTheoremContent[#1]{\GetTheoremDefinitionFormatEnvironmentName}{#2}%
} % End of \newcommand{}

\newcommand{\InitTheoremDefinitionFormat}
{%
  \theoremstyle{definition}%
  \NCKUPrivateInitNumberedTheoremFormat{Definition}{%
    \GetTheoremDefinitionFormatEnvironmentName}{%
    \GetTheoremDefinitionFormatShowText}{%
    \GetTheoremDefinitionFormatFollowCounter}%
} % End of \newcommand{}

% ---------------------------------------------------------

\pgfkeys
{
  /TheoremConditionFormat/.is family, /TheoremConditionFormat,
  default/.style =
  {
    EnvironmentName = {EnvCondition},
    ShowText = {Condition},
    FollowCounter = Section,
  },
  EnvironmentName/.estore in = \GetTheoremConditionFormatEnvironmentName,
  ShowText/.estore in = \GetTheoremConditionFormatShowText,
  FollowCounter/.estore in = \GetTheoremConditionFormatFollowCounter,
} % End of \pgfkeys{}

\newcommand{\InsertCondition}[2][\empty]
{%
  \InsertTheoremContent[#1]{\GetTheoremConditionFormatEnvironmentName}{#2}%
} % End of \newcommand{}

\newcommand{\InitTheoremConditionFormat}
{%
  \theoremstyle{definition}%
  \NCKUPrivateInitNumberedTheoremFormat{Condition}{%
    \GetTheoremConditionFormatEnvironmentName}{%
    \GetTheoremConditionFormatShowText}{%
    \GetTheoremConditionFormatFollowCounter}%
} % End of \newcommand{}

% ---------------------------------------------------------

\pgfkeys
{
  /TheoremProblemFormat/.is family, /TheoremProblemFormat,
  default/.style =
  {
    EnvironmentName = {EnvProblem},
    ShowText = {Problem},
    FollowCounter = Section,
  },
  EnvironmentName/.estore in = \GetTheoremProblemFormatEnvironmentName,
  ShowText/.estore in = \GetTheoremProblemFormatShowText,
  FollowCounter/.estore in = \GetTheoremProblemFormatFollowCounter,
} % End of \pgfkeys{}

\newcommand{\InsertProblem}[2][\empty]
{%
  \InsertTheoremContent[#1]{\GetTheoremProblemFormatEnvironmentName}{#2}%
} % End of \newcommand{}

\newcommand{\InitTheoremProblemFormat}
{%
  \theoremstyle{definition}%
  \NCKUPrivateInitNumberedTheoremFormat{Problem}{%
    \GetTheoremProblemFormatEnvironmentName}{%
    \GetTheoremProblemFormatShowText}{%
    \GetTheoremProblemFormatFollowCounter}%
} % End of \newcommand{}

% ---------------------------------------------------------

\pgfkeys
{
  /TheoremExampleFormat/.is family, /TheoremExampleFormat,
  default/.style =
  {
    EnvironmentName = {EnvExample},
    ShowText = {Example},
    FollowCounter = Section,
  },
  EnvironmentName/.estore in = \GetTheoremExampleFormatEnvironmentName,
  ShowText/.estore in = \GetTheoremExampleFormatShowText,
  FollowCounter/.estore in = \GetTheoremExampleFormatFollowCounter,
} % End of \pgfkeys{}

\newcommand{\InsertExample}[2][\empty]
{%
  \InsertTheoremContent[#1]{\GetTheoremExampleFormatEnvironmentName}{#2}%
} % End of \newcommand{}

\newcommand{\InitTheoremExampleFormat}
{%
  \theoremstyle{definition}%
  \NCKUPrivateInitNumberedTheoremFormat{Example}{%
    \GetTheoremExampleFormatEnvironmentName}{%
    \GetTheoremExampleFormatShowText}{%
    \GetTheoremExampleFormatFollowCounter}%
} % End of \newcommand{}

% ---------------------------------------------------------

\pgfkeys
{
  /TheoremLemmaFormat/.is family, /TheoremLemmaFormat,
  default/.style =
  {
    EnvironmentName = {EnvLemma},
    ShowText = {Lemma},
    FollowCounter = Section,
  },
  EnvironmentName/.estore in = \GetTheoremLemmaFormatEnvironmentName,
  ShowText/.estore in = \GetTheoremLemmaFormatShowText,
  FollowCounter/.estore in = \GetTheoremLemmaFormatFollowCounter,
} % End of \pgfkeys{}

\newcommand{\InsertLemma}[2][\empty]
{%
  \InsertTheoremContent[#1]{\GetTheoremLemmaFormatEnvironmentName}{#2}%
} % End of \newcommand{}

\newcommand{\InitTheoremLemmaFormat}
{%
  \theoremstyle{plain}%
  \NCKUPrivateInitNumberedTheoremFormat{Lemma}{%
    \GetTheoremLemmaFormatEnvironmentName}{%
    \GetTheoremLemmaFormatShowText}{%
    \GetTheoremLemmaFormatFollowCounter}%
} % End of \newcommand{}

% ---------------------------------------------------------

\pgfkeys
{
  /TheoremCorollaryFormat/.is family, /TheoremCorollaryFormat,
  default/.style =
  {
    EnvironmentName = {EnvCorollary},
    ShowText = {Corollary},
    FollowCounter = Section,
  },
  EnvironmentName/.estore in = \GetTheoremCorollaryFormatEnvironmentName,
  ShowText/.estore in = \GetTheoremCorollaryFormatShowText,
  FollowCounter/.estore in = \GetTheoremCorollaryFormatFollowCounter,
} % End of \pgfkeys{}

\newcommand{\InsertCorollary}[2][\empty]
{%
  \InsertTheoremContent[#1]{\GetTheoremCorollaryFormatEnvironmentName}{#2}%
} % End of \newcommand{}

\newcommand{\InitTheoremCorollaryFormat}
{%
  \theoremstyle{plain}%
  \NCKUPrivateInitNumberedTheoremFormat{Corollary}{%
    \GetTheoremCorollaryFormatEnvironmentName}{%
    \GetTheoremCorollaryFormatShowText}{%
    \GetTheoremCorollaryFormatFollowCounter}%
} % End of \newcommand{}

% ---------------------------------------------------------

\pgfkeys
{
  /TheoremPropositionFormat/.is family, /TheoremPropositionFormat,
  default/.style =
  {
    EnvironmentName = {EnvProposition},
    ShowText = {Proposition},
    FollowCounter = Section,
  },
  EnvironmentName/.estore in = \GetTheoremPropositionFormatEnvironmentName,
  ShowText/.estore in = \GetTheoremPropositionFormatShowText,
  FollowCounter/.estore in = \GetTheoremPropositionFormatFollowCounter,
} % End of \pgfkeys{}

\newcommand{\InsertProposition}[2][\empty]
{%
  \InsertTheoremContent[#1]{\GetTheoremPropositionFormatEnvironmentName}{#2}%
} % End of \newcommand{}

\newcommand{\InitTheoremPropositionFormat}
{%
  \theoremstyle{plain}%
  \NCKUPrivateInitNumberedTheoremFormat{Proposition}{%
    \GetTheoremPropositionFormatEnvironmentName}{%
    \GetTheoremPropositionFormatShowText}{%
    \GetTheoremPropositionFormatFollowCounter}%
} % End of \newcommand{}

% ---------------------------------------------------------

\pgfkeys
{
  /TheoremConjectureFormat/.is family, /TheoremConjectureFormat,
  default/.style =
  {
    EnvironmentName = {EnvConjecture},
    ShowText = {Conjecture},
    FollowCounter = Section,
  },
  EnvironmentName/.estore in = \GetTheoremConjectureFormatEnvironmentName,
  ShowText/.estore in = \GetTheoremConjectureFormatShowText,
  FollowCounter/.estore in = \GetTheoremConjectureFormatFollowCounter,
} % End of \pgfkeys{}

\newcommand{\InsertConjecture}[2][\empty]
{%
  \InsertTheoremContent[#1]{\GetTheoremConjectureFormatEnvironmentName}{#2}%
} % End of \newcommand{}

\newcommand{\InitTheoremConjectureFormat}
{%
  \theoremstyle{plain}%
  \NCKUPrivateInitNumberedTheoremFormat{Conjecture}{%
    \GetTheoremConjectureFormatEnvironmentName}{%
    \GetTheoremConjectureFormatShowText}{%
    \GetTheoremConjectureFormatFollowCounter}%
} % End of \newcommand{}

% ---------------------------------------------------------

\pgfkeys
{
  /TheoremHypothesisFormat/.is family, /TheoremHypothesisFormat,
  default/.style =
  {
    EnvironmentName = {EnvHypothesis},
    ShowText = {Hypothesis},
    FollowCounter = Section,
  },
  EnvironmentName/.estore in = \GetTheoremHypothesisFormatEnvironmentName,
  ShowText/.estore in = \GetTheoremHypothesisFormatShowText,
  FollowCounter/.estore in = \GetTheoremHypothesisFormatFollowCounter,
} % End of \pgfkeys{}

\newcommand{\InsertHypothesis}[2][\empty]
{%
  \InsertTheoremContent[#1]{\GetTheoremHypothesisFormatEnvironmentName}{#2}%
} % End of \newcommand{}

\newcommand{\InitTheoremHypothesisFormat}
{%
  \theoremstyle{definition}%
  \NCKUPrivateInitNumberedTheoremFormat{Hypothesis}{%
    \GetTheoremHypothesisFormatEnvironmentName}{%
    \GetTheoremHypothesisFormatShowText}{%
    \GetTheoremHypothesisFormatFollowCounter}%
} % End of \newcommand{}

% ---------------------------------------------------------

\pgfkeys
{
  /TheoremQuestionFormat/.is family, /TheoremQuestionFormat,
  default/.style =
  {
    EnvironmentName = {EnvQuestion},
    ShowText = {Question},
    FollowCounter = Section,
  },
  EnvironmentName/.estore in = \GetTheoremQuestionFormatEnvironmentName,
  ShowText/.estore in = \GetTheoremQuestionFormatShowText,
  FollowCounter/.estore in = \GetTheoremQuestionFormatFollowCounter,
} % End of \pgfkeys{}

\newcommand{\InsertQuestion}[2][\empty]
{%
  \InsertTheoremContent[#1]{\GetTheoremQuestionFormatEnvironmentName}{#2}%
} % End of \newcommand{}

\newcommand{\InitTheoremQuestionFormat}
{%
  \theoremstyle{definition}%
  \NCKUPrivateInitNumberedTheoremFormat{Question}{%
    \GetTheoremQuestionFormatEnvironmentName}{%
    \GetTheoremQuestionFormatShowText}{%
    \GetTheoremQuestionFormatFollowCounter}%
} % End of \newcommand{}

% ---------------------------------------------------------

\pgfkeys
{
  /TheoremAssertionFormat/.is family, /TheoremAssertionFormat,
  default/.style =
  {
    EnvironmentName = {EnvAssertion},
    ShowText = {Assertion},
    FollowCounter = Section,
  },
  EnvironmentName/.estore in = \GetTheoremAssertionFormatEnvironmentName,
  ShowText/.estore in = \GetTheoremAssertionFormatShowText,
  FollowCounter/.estore in = \GetTheoremAssertionFormatFollowCounter,
} % End of \pgfkeys{}

\newcommand{\InsertAssertion}[2][\empty]
{%
  \InsertTheoremContent[#1]{\GetTheoremAssertionFormatEnvironmentName}{#2}%
} % End of \newcommand{}

\newcommand{\InitTheoremAssertionFormat}
{%
  \theoremstyle{plain}%
  \NCKUPrivateInitNumberedTheoremFormat{Assertion}{%
    \GetTheoremAssertionFormatEnvironmentName}{%
    \GetTheoremAssertionFormatShowText}{%
    \GetTheoremAssertionFormatFollowCounter}%
} % End of \newcommand{}

% ---------------------------------------------------------

\pgfkeys
{
  /TheoremCriterionFormat/.is family, /TheoremCriterionFormat,
  default/.style =
  {
    EnvironmentName = {EnvCriterion},
    ShowText = {Criterion},
    FollowCounter = Section,
  },
  EnvironmentName/.estore in = \GetTheoremCriterionFormatEnvironmentName,
  ShowText/.estore in = \GetTheoremCriterionFormatShowText,
  FollowCounter/.estore in = \GetTheoremCriterionFormatFollowCounter,
} % End of \pgfkeys{}

\newcommand{\InsertCriterion}[2][\empty]
{%
  \InsertTheoremContent[#1]{\GetTheoremCriterionFormatEnvironmentName}{#2}%
} % End of \newcommand{}

\newcommand{\InitTheoremCriterionFormat}
{%
  \theoremstyle{plain}%
  \NCKUPrivateInitNumberedTheoremFormat{Criterion}{%
    \GetTheoremCriterionFormatEnvironmentName}{%
    \GetTheoremCriterionFormatShowText}{%
    \GetTheoremCriterionFormatFollowCounter}%
} % End of \newcommand{}

% ---------------------------------------------------------

\pgfkeys
{
  /TheoremProofFormat/.is family, /TheoremProofFormat,
  default/.style =
  {
    EnvironmentName = {EnvProof},
    ShowText = {Proof},
    FollowCounter = \empty,
  },
  EnvironmentName/.estore in = \GetTheoremProofFormatEnvironmentName,
  ShowText/.estore in = \GetTheoremProofFormatShowText,
  FollowCounter/.estore in = \GetTheoremProofFormatFollowCounter,
} % End of \pgfkeys{}

\newcommand{\InsertProof}[1]
{%
  \InsertTheoremContent[\empty]{%
    \GetTheoremProofFormatEnvironmentName}{#1}%
} % End of \newcommand{}

\newcommand{\InitTheoremProofFormat}
{%
  \theoremstyle{definition}%
  \ifthenelse{\equal{\GetTheoremProofFormatFollowCounter}{\empty}}%
  {%
    \newtheorem*{%
      \GetTheoremProofFormatEnvironmentName}{%
      \GetTheoremProofFormatShowText}
  }%
  {%
    %
    \MappingTheoremCounter[Proof]{\GetTheoremProofFormatFollowCounter}%
    %
    \newtheorem{%
      \GetTheoremProofFormatEnvironmentName}{%
      \GetTheoremProofFormatShowText}[%
      \GetTheoremProofFormatFollowCounter]%
  }%
} % End of \newcommand{}

% ---------------------------------------------------------

\pgfkeys
{
  /TheoremNoteFormat/.is family, /TheoremNoteFormat,
  default/.style =
  {
    EnvironmentName = {EnvNote},
    ShowText = {Note},
    FollowCounter = \empty,
  },
  EnvironmentName/.estore in = \GetTheoremNoteFormatEnvironmentName,
  ShowText/.estore in = \GetTheoremNoteFormatShowText,
  FollowCounter/.estore in = \GetTheoremNoteFormatFollowCounter,
} % End of \pgfkeys{}

\newcommand{\InsertNote}[1]
{%
  \InsertTheoremContent[\empty]{%
    \GetTheoremNoteFormatEnvironmentName}{#1}%
} % End of \newcommand{}

\newcommand{\InitTheoremNoteFormat}
{%
  \theoremstyle{definition}%
  \ifthenelse{\equal{\GetTheoremNoteFormatFollowCounter}{\empty}}%
  {%
    \newtheorem*{%
      \GetTheoremNoteFormatEnvironmentName}{%
      \GetTheoremNoteFormatShowText}
  }%
  {%
    %
    \MappingTheoremCounter[Note]{\GetTheoremNoteFormatFollowCounter}%
    %
    \newtheorem{%
      \GetTheoremNoteFormatEnvironmentName}{%
      \GetTheoremNoteFormatShowText}[%
      \GetTheoremNoteFormatFollowCounter]%
  }%
} % End of \newcommand{}

% ---------------------------------------------------------

\pgfkeys
{
  /TheoremAnnotationFormat/.is family, /TheoremAnnotationFormat,
  default/.style =
  {
    EnvironmentName = {EnvAnnotation},
    ShowText = {Annotation},
    FollowCounter = \empty,
  },
  EnvironmentName/.estore in = \GetTheoremAnnotationFormatEnvironmentName,
  ShowText/.estore in = \GetTheoremAnnotationFormatShowText,
  FollowCounter/.estore in = \GetTheoremAnnotationFormatFollowCounter,
} % End of \pgfkeys{}

\newcommand{\InsertAnnotation}[1]
{%
  \InsertTheoremContent[\empty]{%
    \GetTheoremAnnotationFormatEnvironmentName}{#1}%
} % End of \newcommand{}

\newcommand{\InitTheoremAnnotationFormat}
{%
  \theoremstyle{definition}%
  \ifthenelse{\equal{\GetTheoremAnnotationFormatFollowCounter}{\empty}}%
  {%
    \newtheorem*{%
      \GetTheoremAnnotationFormatEnvironmentName}{%
      \GetTheoremAnnotationFormatShowText}
  }%
  {%
    %
    \MappingTheoremCounter[Annotation]{\GetTheoremAnnotationFormatFollowCounter}%
    %
    \newtheorem{%
      \GetTheoremAnnotationFormatEnvironmentName}{%
      \GetTheoremAnnotationFormatShowText}[%
      \GetTheoremAnnotationFormatFollowCounter]%
  }%
} % End of \newcommand{}

% ---------------------------------------------------------

\pgfkeys
{
  /TheoremClaimFormat/.is family, /TheoremClaimFormat,
  default/.style =
  {
    EnvironmentName = {EnvClaim},
    ShowText = {Claim},
    FollowCounter = \empty,
  },
  EnvironmentName/.estore in = \GetTheoremClaimFormatEnvironmentName,
  ShowText/.estore in = \GetTheoremClaimFormatShowText,
  FollowCounter/.estore in = \GetTheoremClaimFormatFollowCounter,
} % End of \pgfkeys{}

\newcommand{\InsertClaim}[1]
{%
  \InsertTheoremContent[\empty]{%
    \GetTheoremClaimFormatEnvironmentName}{#1}%
} % End of \newcommand{}

\newcommand{\InitTheoremClaimFormat}
{%
  \theoremstyle{definition}%
  \ifthenelse{\equal{\GetTheoremClaimFormatFollowCounter}{\empty}}%
  {%
    \newtheorem*{%
      \GetTheoremClaimFormatEnvironmentName}{%
      \GetTheoremClaimFormatShowText}
  }%
  {%
    %
    \MappingTheoremCounter[Claim]{\GetTheoremClaimFormatFollowCounter}%
    %
    \newtheorem{%
      \GetTheoremClaimFormatEnvironmentName}{%
      \GetTheoremClaimFormatShowText}[%
      \GetTheoremClaimFormatFollowCounter]%
  }%
} % End of \newcommand{}

% ---------------------------------------------------------

\pgfkeys
{
  /TheoremCaseFormat/.is family, /TheoremCaseFormat,
  default/.style =
  {
    EnvironmentName = {EnvCase},
    ShowText = {Case},
    FollowCounter = \empty,
  },
  EnvironmentName/.estore in = \GetTheoremCaseFormatEnvironmentName,
  ShowText/.estore in = \GetTheoremCaseFormatShowText,
  FollowCounter/.estore in = \GetTheoremCaseFormatFollowCounter,
} % End of \pgfkeys{}

\newcommand{\InsertCase}[1]
{%
  \InsertTheoremContent[\empty]{%
    \GetTheoremCaseFormatEnvironmentName}{#1}%
} % End of \newcommand{}

\newcommand{\InitTheoremCaseFormat}
{%
  \theoremstyle{definition}%
  \ifthenelse{\equal{\GetTheoremCaseFormatFollowCounter}{\empty}}%
  {%
    \newtheorem*{%
      \GetTheoremCaseFormatEnvironmentName}{%
      \GetTheoremCaseFormatShowText}
  }%
  {%
    %
    \MappingTheoremCounter[Case]{\GetTheoremCaseFormatFollowCounter}%
    %
    \newtheorem{%
      \GetTheoremCaseFormatEnvironmentName}{%
      \GetTheoremCaseFormatShowText}[%
      \GetTheoremCaseFormatFollowCounter]%
  }%
} % End of \newcommand{}

% ---------------------------------------------------------

\pgfkeys
{
  /TheoremAcknowledgmentFormat/.is family, /TheoremAcknowledgmentFormat,
  default/.style =
  {
    EnvironmentName = {EnvAcknowledgment},
    ShowText = {Acknowledgment},
    FollowCounter = \empty,
  },
  EnvironmentName/.estore in = \GetTheoremAcknowledgmentFormatEnvironmentName,
  ShowText/.estore in = \GetTheoremAcknowledgmentFormatShowText,
  FollowCounter/.estore in = \GetTheoremAcknowledgmentFormatFollowCounter,
} % End of \pgfkeys{}

\newcommand{\InsertAcknowledgment}[1]
{%
  \InsertTheoremContent[\empty]{%
    \GetTheoremAcknowledgmentFormatEnvironmentName}{#1}%
} % End of \newcommand{}

\newcommand{\InitTheoremAcknowledgmentFormat}
{%
  \theoremstyle{definition}%
  \ifthenelse{\equal{\GetTheoremAcknowledgmentFormatFollowCounter}{\empty}}%
  {%
    \newtheorem*{%
      \GetTheoremAcknowledgmentFormatEnvironmentName}{%
      \GetTheoremAcknowledgmentFormatShowText}
  }%
  {%
    %
    \MappingTheoremCounter[Acknowledgment]{\GetTheoremAcknowledgmentFormatFollowCounter}%
    %
    \newtheorem{%
      \GetTheoremAcknowledgmentFormatEnvironmentName}{%
      \GetTheoremAcknowledgmentFormatShowText}[%
      \GetTheoremAcknowledgmentFormatFollowCounter]%
  }%
} % End of \newcommand{}

% ---------------------------------------------------------

\pgfkeys
{
  /TheoremConclusionFormat/.is family, /TheoremConclusionFormat,
  default/.style =
  {
    EnvironmentName = {EnvConclusion},
    ShowText = {Conclusion},
    FollowCounter = \empty,
  },
  EnvironmentName/.estore in = \GetTheoremConclusionFormatEnvironmentName,
  ShowText/.estore in = \GetTheoremConclusionFormatShowText,
  FollowCounter/.estore in = \GetTheoremConclusionFormatFollowCounter,
} % End of \pgfkeys{}

\newcommand{\InsertConclusion}[1]
{%
  \InsertTheoremContent[\empty]{%
    \GetTheoremConclusionFormatEnvironmentName}{#1}%
} % End of \newcommand{}

\newcommand{\InitTheoremConclusionFormat}
{%
  \theoremstyle{definition}%
  \ifthenelse{\equal{\GetTheoremConclusionFormatFollowCounter}{\empty}}%
  {%
    \newtheorem*{%
      \GetTheoremConclusionFormatEnvironmentName}{%
      \GetTheoremConclusionFormatShowText}
  }%
  {%
    %
    \MappingTheoremCounter[Conclusion]{\GetTheoremConclusionFormatFollowCounter}%
    %
    \newtheorem{%
      \GetTheoremConclusionFormatEnvironmentName}{%
      \GetTheoremConclusionFormatShowText}[%
      \GetTheoremConclusionFormatFollowCounter]%
  }%
} % End of \newcommand{}

% ---------------------------------------------------------

\pgfkeys
{
  /TheoremSummaryFormat/.is family, /TheoremSummaryFormat,
  default/.style =
  {
    EnvironmentName = {EnvSummary},
    ShowText = {Summary},
    FollowCounter = \empty,
  },
  EnvironmentName/.estore in = \GetTheoremSummaryFormatEnvironmentName,
  ShowText/.estore in = \GetTheoremSummaryFormatShowText,
  FollowCounter/.estore in = \GetTheoremSummaryFormatFollowCounter,
} % End of \pgfkeys{}

\newcommand{\InsertSummary}[1]
{%
  \InsertTheoremContent[\empty]{%
    \GetTheoremSummaryFormatEnvironmentName}{#1}%
} % End of \newcommand{}

\newcommand{\InitTheoremSummaryFormat}
{%
  \theoremstyle{definition}%
  \ifthenelse{\equal{\GetTheoremSummaryFormatFollowCounter}{\empty}}%
  {%
    \newtheorem*{%
      \GetTheoremSummaryFormatEnvironmentName}{%
      \GetTheoremSummaryFormatShowText}
  }%
  {%
    %
    \MappingTheoremCounter[Summary]{\GetTheoremSummaryFormatFollowCounter}%
    %
    \newtheorem{%
      \GetTheoremSummaryFormatEnvironmentName}{%
      \GetTheoremSummaryFormatShowText}[%
      \GetTheoremSummaryFormatFollowCounter]%
  }%
} % End of \newcommand{}

% ---------------------------------------------------------

% Private ordered registry. Keep per-type wrappers, key families, styles, and
% counter implementations independent until a later fixture-backed slice.
\newcommand{\NCKUPrivateForEachTheoremFormat}[1]
{%
  #1{Definition}%
  #1{Condition}%
  #1{Problem}%
  #1{Example}%
  #1{Theorem}%
  #1{Lemma}%
  #1{Corollary}%
  #1{Proposition}%
  #1{Conjecture}%
  #1{Proof}%
  #1{Note}%
  #1{Annotation}%
  #1{Claim}%
  #1{Case}%
  #1{Acknowledgment}%
  #1{Conclusion}%
  #1{Criterion}%
  #1{Assertion}%
  #1{Question}%
  #1{Hypothesis}%
  #1{Summary}%
} % End of \newcommand{}

\newcommand{\NCKUPrivateRegisterTheoremFormat}[1]
{%
  \expandafter\def\csname NCKUPrivateTheoremFormatRegistered@#1\endcsname{}%
} % End of \newcommand{}

\NCKUPrivateForEachTheoremFormat{\NCKUPrivateRegisterTheoremFormat}

\newcommand{\NCKUPrivateInitTheoremFormat}[1]
{%
  \csname InitTheorem#1Format\endcsname%
} % End of \newcommand{}

\newcommand{\InitTheoremFormats}
{%
  \NCKUPrivateForEachTheoremFormat{\NCKUPrivateInitTheoremFormat}%
} % End of \newcommand{}

\newcommand{\SetTheoremFormat}[2][\empty]
{%
  \ifcsname NCKUPrivateTheoremFormatRegistered@#1\endcsname
    \pgfkeys{/Theorem#1Format, default, #2}%
  \fi
} % End of \newcommand{}

% ---------------------------------------------------------

\SetTheoremFormat[Definition]{ShowText = {Definition}}%
\SetTheoremFormat[Condition]{ShowText = {Condition}}%
\SetTheoremFormat[Problem]{ShowText = {Problem}}%
\SetTheoremFormat[Example]{ShowText = {Example}}%
\SetTheoremFormat[Theorem]{ShowText = {Theorem}}%
\SetTheoremFormat[Lemma]{ShowText = {Lemma}}%
\SetTheoremFormat[Corollary]{ShowText = {Corollary}}%
\SetTheoremFormat[Proposition]{ShowText = {Proposition}}%
\SetTheoremFormat[Conjecture]{ShowText = {Conjecture}}%
\SetTheoremFormat[Proof]{ShowText = {Proof}}%
\SetTheoremFormat[Note]{ShowText = {Note}}%
\SetTheoremFormat[Annotation]{ShowText = {Annotation}}%
\SetTheoremFormat[Claim]{ShowText = {Claim}}%
\SetTheoremFormat[Case]{ShowText = {Case}}%
\SetTheoremFormat[Acknowledgment]{ShowText = {Acknowledgment}}%
\SetTheoremFormat[Conclusion]{ShowText = {Conclusion}}%
\SetTheoremFormat[Criterion]{ShowText = {Criterion}}%
\SetTheoremFormat[Assertion]{ShowText = {Assertion}}%
\SetTheoremFormat[Question]{ShowText = {Question}}%
\SetTheoremFormat[Hypothesis]{ShowText = {Hypothesis}}%
\SetTheoremFormat[Summary]{ShowText = {Summary}}%

% ----------------------------------------------------------------------------
